% =============================================================================
% §8 R4 — THE CEILING  ★ write this section FIRST, with §3 (ke-hoach §4)
%
% 🔴🔴 RED LINE A2 (prereg §6/A2 · CLAUDE.md lằn ranh 5) — READ BEFORE WRITING:
%   1. "We establish the ceiling" / "we are the first to show a ceiling" is
%      ABSOLUTELY FORBIDDEN. François & Tinarrage (JMIV 2026) PUBLISHED the
%      concept + the dataset (BraTS FLAIR) + the number (0.83±0.18) + the DL
%      comparison. Overclaiming here, in a paper that punishes overclaiming,
%      is death with no defense.
%   2. The level-set / purity-prefix theorem is NOT OURS. It belongs to Lipton,
%      Elkan & Narayanaswamy (arXiv:1402.1892) and RankSEG. Dice = F1; purity is
%      a calibrated posterior. Shipping it as "Theorem 1 (ours)" repeats the
%      exact fabricated-novelty offense this paper exists to document.
%      => CLAIM THE APPLICATION AND THE DECOMPOSITION. NEVER THE MATHEMATICS.
%   3. The old statement "the single-interval oracle is the ceiling of ALL
%      intensity thresholding" is MATHEMATICALLY FALSE — a non-contiguous
%      gray-level set beats it. The correct tightest bound is the arbitrary
%      gray-level-set oracle.
%
% 🔴 PREREGISTERED FALLBACK (locked BEFORE seeing numbers, prereg §6/A2): the
% oracle bound on WT/FLAIR is NOT low (0.83 sits inside the inter-rater
% consensus band), so a FLAIR-only 2D U-Net may well NOT beat it. If it does not,
% the headline does NOT retreat to "thresholding saturates" (that is exactly what
% François printed = zero novelty). It becomes: "the ceiling is HIGH — so the
% failure of thresholding is not a limit of the representation but of the
% THRESHOLD-SELECTION problem, which no metaheuristic in this line addresses."
%
% Oracles are NOT methods (A3, family C): label "uses test-time ground truth" on
% EVERY row and EVERY figure element.
% NO RESULTS PROSE until results/ceiling/ and results/unet/ exist.
% =============================================================================
\section{Decomposing the Ceiling}
\label{sec:ceiling}

\subsection{What Is Already Known}

The existence of an upper bound on thresholding-based segmentation of brain MRI
is an established result, reported on BraTS FLAIR by François and
Tinarrage~\cite{francois_jmiv}. We do not restate it as a finding of ours. Our
question is different: \emph{what does that bound consist of?}

\subsection{The Tightest Intensity-Only Bound}

% APPLICATION, not mathematics. Attribute in the same breath as the statement.
For a fixed image, the mask that maximizes Dice among all masks defined by an
arbitrary set of gray levels is obtained by ordering gray levels by the purity
of their voxels with respect to the reference mask and taking the best prefix;
this is the F1-threshold result of Lipton \emph{et al.}~\cite{lipton_f1}
transposed to the segmentation setting, where Dice coincides with F1 and purity
plays the role of a calibrated posterior. We claim no mathematical novelty for
it. Its consequence here is practical: the bound is computable in $O(L \log L)$
from a histogram, so the attainable performance of the entire thresholding
family is knowable before any optimizer is written.

\TODO{state the bound formally, with attribution in the statement itself, and
note that it is strictly tighter than the contiguous-band oracle because a
non-contiguous gray-level set is admissible. Do NOT label it Theorem 1 (ours).}

\subsection{The Ladder}

\TODO{describe the ladder on a single Dice axis over the same held-out patients:
random search at the shared budget; the metaheuristics; the exact optimum;
classical thresholds and clustering baselines; the out-of-fold single-parameter
threshold; the contiguous-band oracle; the arbitrary-gray-level-set oracle; a
2D U-Net on identical single-slice input; and, as CONTEXT ONLY, a literature
value for a 3D multi-modal network~\cite{isensee_nnunet}, explicitly labeled
``reference from literature, not run by us, different input protocol'' and never
compared as a baseline.}

% --- FIGURE 4 — ceiling ladder ★ the most valuable figure in the paper -------
% results/ceiling/ + results/unet/ <- scripts/build_figures.py
%   -> figures/fig4_ceiling_ladder.pdf
\begin{figure}[t]
  \centering
  \inputfigure{fig4_ceiling_ladder}
  \caption{Ceiling ladder on a common Dice axis over the same held-out patients.
  \TODO{caption: every oracle rung labeled ``uses test-time ground truth'';
  n patients, folds, seeds; the literature reference rung labeled as not run by
  us and using a different input protocol.}}
  \label{fig:ladder}
\end{figure}

\subsection{Decomposition}

\TODO{the two gaps that constitute the contribution: (i) band oracle to
gray-level-set oracle = what is lost by insisting on a contiguous intensity
band; (ii) gray-level-set oracle to the learned model on identical input = what
is absent from intensity but present in the pixels. Report both with CIs, from
results/. This decomposition, not the existence of the ceiling, is what is ours.}

\subsection{Where the Representational Limit Is Provably Real}

\TODO{the hard target: enhancing tumor on T1ce is a bright rim around a dark
core and therefore cannot be represented by any single contiguous intensity
band; the favorable target (whole tumor on FLAIR) is retained as the best case
for the methods under audit, making the argument a fortiori.}

\TODO{preregistered branch statement: if the learned model does not exceed the
oracle bound, report that plainly and adopt the preregistered headline (the
ceiling is high; the failure is one of threshold SELECTION, not of
representation). Do not retreat to a restatement of prior work.}
